\documentclass[a4paper,12pt]{article}

\usepackage[left=2cm, top=2cm, right=2cm, bottom=1cm, includeheadfoot,
    headheight=50pt]{geometry}
\usepackage{fancyhdr}
\usepackage{lastpage}
\usepackage{float}
\usepackage[most]{tcolorbox}
\usepackage{enumitem}
\usepackage{graphicx}

\usepackage{amsmath, amsthm}

\usepackage{fontspec}
\usepackage{unicode-math}

\setmainfont{Linux Libertine O}

\newtheorem{theorem}{Theorem}[section]
\newtheorem{lemma}{Lemma}[section]

\newcommand{\Name}{Hostek}
\newcommand{\Email}{your.email@example.com}
\newcommand{\ProblemNumber}{LXXV OF, etap 2, zadanie 1}

\pagestyle{fancy}
\fancyhf{}
\fancyhead[L]{\Name \\ \Email}
\fancyhead[C]{\ProblemNumber}
\fancyfoot[C]{\thepage/\pageref{LastPage}}

\renewcommand{\headrulewidth}{0.4pt}
\renewcommand{\footrulewidth}{0.4pt}

\begin{document}

\fontsize{15}{18}\selectfont

\section*{Problem Statement}

Baron Münchhausen postanowił wznieść się jak najwyżej wykorzystując armatkę wystrzeliwującą pionowo w górę pocisk o masie $m$ z prędkością $v_0$. Baron trzyma w rękach kołowrotek z częściowo nawiniętą nierozciągliwą, wiotką, nieważką liną, której drugi koniec jest przymocowany do pocisku. Kołowrotek ma wbudowany hamulec powodujący, że w trakcie rozwijania liny jej naprężenie wynosi $N_0$ (gdy naprężenie jest mniejsze od $N_0$ lina się nie rozwija, a długość rozwiniętej części pozostaje stała). Początkowa długość rozwiniętej części liny wynosi $l_0$. Barona potraktuj jako bryłę sztywną o masie $M$.

Na jaką maksymalną wysokość wzniesie się baron, zakładając, że całkowita długość liny jest wystarczająco duża. Dodatkowo wyznacz ciepło, jakie wydzieli się na kołowrotku w trakcie wznoszenia się barona.

Pomiń rozmiary liniowe barona, armatki i pocisku oraz opór powietrza. Przyspieszenie ziemskie wynosi $g$. Początkowo baron znajduje się tuż obok armatki; zakładamy, że pocisk nie uderza w niego.

Podaj wyniki liczbowe dla:
$m = 10 \text{ kg}$, $M = 80 \text{ kg}$, $l_0 = 5 \text{ m}$, $v_0 = 200 \text{ m/s}$, $N_0 = 8000 \text{ N}$, $g = 9,81 \text{ m/s}^2$.

\bigskip

\subsection*{Solution:}

\end{document}